\section{Experimental Setup}
\label{sec:setup}

\subsection{Datasets}

\emph{CUHK-CQRC}~\cite{yang2018hiq} contains photographs of printed three-layer HiQ symbols
at versions 21, 25, 30, 35 and 40, taken with several phones under different lighting. Our
copy holds the \HiqAttempted{} images of the public release in eight capture folders. The
release's metadata marks \SplitTest{} of them as HiQ's test set, \SplitTrain{} as the images
HiQ's models were trained on, and \SplitUnused{} as unused; the HiQ paper reports 5390 test
samples~\cite{yang2018hiq}. Our method trains nothing, so its results use every image; the
comparison with the HiQ decoder uses the test set alone. One image appears twice under two
names. Ground truth for each image is the clean source render of its version.

The \emph{qrism dataset} is new and will be released with this paper. It contains
\QrAttempted{} photographs of qrism symbols at versions 1, 5, 10, 15, 25, 30 and 40, each
at all four error-correction levels. Modules measure 0.5 and 0.84\,mm, or 0.34\,mm in the
small-module condition, and each symbol was photographed five times under the
conditions of Table~\ref{tab:dataset} with four phones: Apple iPhone~13 and iPhone~16,
Motorola Edge~50 Fusion and OnePlus Nord~CE6. Symbols were printed at 600\,dpi on plain A4 paper with an HP laser printer, from renders of
20, 12 and 8 pixels per module. The on-screen condition used an LG display, zoomed so that
modules have the same sizes as in print. Photographs were taken about 20 to 25\,cm from the symbol. The
0.34\,mm symbols, and the 0.84\,mm symbols of the close condition, were taken from about 10 to
15\,cm.

% Generated by scripts/make_tables.py. Do not edit.
\begin{table}[t]
\centering
\small
\caption{The qrism dataset by capture condition: phones used, module sizes of the
source renders in pixels, photos, and photos the localizer found at the right version.}
\label{tab:dataset}
\begin{tabular}{@{}lrlrr@{}}
\toprule
Condition & Phones & px & Photos & Localized \\
\midrule
Reference capture & 4 & 12, 20 & 1120 & 1024 \\
Outdoor daylight & 1 & 12, 20 & 280 & 271 \\
Phone flash on & 2 & 12, 20 & 560 & 458 \\
Shown on a display & 1 & 12, 20 & 280 & 75 \\
Oblique view & 1 & 12, 20 & 40 & 0 \\
Close range & 1 & 20 & 28 & 28 \\
8 px modules & 2 & 8 & 280 & 259 \\
White balance cool & 1 & 12, 20 & 280 & 23 \\
White balance neutral & 1 & 12, 20 & 280 & 145 \\
White balance warm & 1 & 12, 20 & 280 & 173 \\
\midrule
Total & 4 & & 3428 & 2456 \\
\bottomrule
\end{tabular}
\end{table}


\subsection{Protocol}

Every color recovery pipeline receives identical input: the same localization and the
same one sample per module, from the library's localizer (Section~\ref{sec:symbol}). Only
stage~(3) of the pipeline differs between rows of a table. A photograph that the localizer
misses, or reads at a version with no source render, is excluded from every color
recovery rate and counted separately. On CUHK-CQRC \HiqLocalized{} of \HiqAttempted{}
photographs are localized (\HiqNoSymbol{} not found, \HiqMisread{} read at a wrong version);
on the qrism dataset \QrLocalized{} of \QrAttempted{} are (\QrNotLocalized{} not found or
read at a wrong version). All rates on a dataset therefore share one denominator, which
lets two settings be compared photograph by photograph.

A predicted HiQ grid is split into its three layers, each layer is rendered as a
monochrome symbol and decoded with the library's standard QR reader, and the symbol counts
as decoded when all three messages equal those of the source. For a predicted qrism grid
the format information must match the source; the three channels are then decoded with
their Reed-Solomon blocks and the message compared with the source.

\subsection{Metrics and statistics}

\emph{Module accuracy} is the share of non-function modules whose color is right; function
modules are left out because they overlap the calibration set. The \emph{bit error rate}
(BER) is the share of layer bits that are wrong before Reed-Solomon correction, over the
same modules, as reported by HiQ~\cite{yang2018hiq}. The \emph{layer decode rate} counts
HiQ layers whose message is right. The \emph{symbol decode rate} counts symbols decoded in
full; HiQ's decoding failure rate is 100 minus this rate.

Decode rates carry 95\% Wilson intervals~\cite{wilson1927}. Two settings on the same
photographs are compared with an exact McNemar test~\cite{mcnemar1947} on the photographs
that one decodes and the other does not. The two parameters, $\beta$ and $\lambda$, are
chosen by symbol decode rate on each dataset separately, and we report whether the choice
made on one dataset holds on the other.

\subsection{Baselines}

The color recovery baselines all use the same calibration modules as our method.
\emph{Nearest centroid} assigns the closest calibration median after a per-channel stretch
that maps the black and white medians to 0 and~1. \emph{Interference cancellation} follows
Blasinski et al.~\cite{blasinski2013pccc}: a non-negative $3\times 3$ mixing matrix is
fitted to the calibration samples in optical density and inverted, and each recovered
channel is thresholded midway between the nearest on and off calibration colors. We also run
the hard-decision classifier of JAB Code~\cite{iso23634} with its image balancing, using one
palette of calibration medians because no finder in either dataset carries all eight colors. \emph{RDA} without leak inversion ($\beta=0$) measures what leak inversion
adds. As an alternative to the fixed $\beta$, we estimate a $\beta$ per channel and
photograph from the width of the blur across the finder edges.

HiQ~\cite{yang2018hiq} is compared in two ways, both with its shipped models: QDA, QDA
with cross-module interference correction (CMI), and a layered SVM. End to end, we run the
released HiQ decoder, which uses its own localization and sampling. For a classifier-only
comparison, we port the three models unchanged to our benchmark and run them on our module
samples after HiQ's own white-point normalization, so that only the classifier differs from
our method. The released CMI has two defects. Its weighting step keeps only the center
module's term and drops the four neighbor terms, so QDA with and without CMI give identical
results; we report a corrected version that applies all five weights. Its sampler also
gathers only the left and upper neighbors for most modules and leaves the other two at zero.
The end-to-end corrected CMI keeps this sampler; on our samples we report it both ways
(Section~\ref{sec:hiq}).
